\chapter{Variables and Control}\label{ch:control}

This chapter covers the forms a program uses to name values and to direct
the flow of evaluation: defining and binding variables, assigning them,
choosing between alternatives, repeating, returning several values, and
leaving a computation early. They are the subject of
\dylanbook{chapters 4, 7 and 8}. Most of them are syntax forms, special
forms or macros (Section~\ref{sec:syntax-forms}); for each, this chapter says which
of its parts are evaluated, in what order, and what it returns.

\section{Module variables: define}\label{sec:control-define}

\index{define@\code{define}}
\begin{quote}
\code{(define \meta{variable-name} \meta{initial-value})}
\end{quote}

\noindent \code{define} creates a module variable named
\meta{variable-name} in the current module and stores in it the value of
\meta{initial-value}. The name is not evaluated; it must be a symbol or a
setter name \code{(setter \meta{symbol})}. The initial value is evaluated,
and if it returns several values, the first is used. At the interactive
prompt \code{define} returns the variable's name.

\begin{dylan}
(define foo 10)
(define bar foo)
(+ foo bar)               ; => 20
(define (setter my-place) 0)
(setter my-place)         ; => 0
(define v (values 1 2))
v                         ; => 1
\end{dylan}

\noindent A module variable can be read from anywhere in its module, by
code before or after its definition, though it has no value until the
definition has been evaluated (Section~\ref{sec:syntax-when}). A variable created by
\code{define} may be assigned with \code{set!}
(Section~\ref{sec:control-assignment}).

Defining a variable that already exists in the module is allowed by the
Common Lisp implementation, which warns \code{redefining variable} and
stores the new value; in the compiler to C, defining a variable twice with
\code{define} simply stores twice, but defining a name both with
\code{define} and as a class or generic function is an error when
compiling.

\begin{notebox}
Later design notes changed \code{define}: one renamed it
\code{define-constant} and added \code{define-variable}, another let it
declare types and receive multiple values, as \code{bind} does. None of
these is part of this language. There is no \code{define-variable} or
\code{define-constant}, and \code{define} takes exactly one plain name and
one expression; \code{(define (x <integer>) 3)} and \code{(define a b
(values 1 2))} are errors in both implementations. The compiler to C adds
\code{define-thread-variable}, which defines a variable each thread has its
own value of (Chapter~\ref{ch:extensions}).
\end{notebox}

\section{Lexical variables: bind}\label{sec:control-bind}

\index{bind@\code{bind}}
\begin{quote}
\code{(bind (\meta{binding} \ldots) \meta{form} \ldots)}
\end{quote}

\noindent \code{bind} creates lexical variables and evaluates the
\meta{form}s, its \term{body}, with them in scope. It returns the values of
the last form, or \code{\#f} if there are no forms. Each \meta{binding} is
a list of one or more variables followed by an expression, the
\meta{init}:

\begin{quote}
\code{(\meta{variable} \ldots{} [\#rest \meta{rest-variable}] \meta{init})}
\end{quote}

\noindent where each \meta{variable} is a name, or a list \code{(\meta{name}
\meta{type})} that declares the variable's type.

The bindings are made one after another. Each init is evaluated in an
environment that includes the variables of the bindings before it, so a
later init can use an earlier variable:

\begin{dylan}
(bind ((x 20)
       (y (+ x x)))
  (+ y y))                ; => 80
(bind ((x 1)))            ; => #f
\end{dylan}

\noindent A variable made by \code{bind} shadows any module variable or
outer lexical variable of the same name, for the extent of the body:

\begin{dylan}
(define foo 10)
(bind ((foo 35))
  (+ foo foo))            ; => 70
(bind ((foo 20))
  (bind ((foo 50))
    (+ foo foo)))         ; => 100
foo                       ; => 10
\end{dylan}

\subsection{Receiving multiple values}

When a binding has several variables, they receive the values of its init
in order. Missing values are \code{\#f}, and values left over are ignored,
unless the variables end with \code{\#rest} and a name, in which case that
variable receives a sequence of the remaining values.

%! print
\begin{dylan}
(bind ((a b c (values 1 2 3)))
  (list a b c))           ; => (1 2 3)
(bind ((a b (values 1)))
  (list a b))             ; => (1 #f)
(bind ((first #rest others (values 1 2 3)))
  others)                 ; => (2 3)
(bind ((#rest all (values)))
  (empty? all))           ; => #t
\end{dylan}

\noindent Multiple values also give parallel binding, like \code{let} in
Lisp, since the init is evaluated before any of its variables exist:

%! print
\begin{dylan}
(bind ((x 10) (y 20))
  (bind ((x y (values y x)))
    (list x y)))          ; => (20 10)
\end{dylan}

\subsection{Type declarations}\label{sec:control-bind-types}

A variable written \code{(\meta{name} \meta{type})} must receive an
instance of \meta{type}, or an error is signaled. The type is an
expression, evaluated before the binding's init; it must evaluate to a class
or a singleton (Section~\ref{sec:functions-singletons}).

\begin{dylan}
(bind (((n <integer>) (+ 1 2)))
  n)                      ; => 3
(bind (((x <integer>) (sqrt 2)))
  x)                      ; => error: 1.4142135623730951 is not
                          ;    an instance of <integer>
\end{dylan}

\noindent The type is checked when the variable is bound. It is not checked
again when the variable is assigned: in both implementations \code{(bind
(((x <integer>) 1)) (set! x 'a) x)} returns \code{a}. (A later design note
made the restriction apply to assignment too; the 1992 manual does not.)
The compiler to C uses declared types to choose faster code
(Chapter~\ref{ch:performance}).

\section{Assignment: set!}\label{sec:control-assignment}

\index{set"!@\code{set"!}}
\begin{quote}
\code{(set! \meta{place} \meta{new-value})}
\end{quote}

\noindent \code{set!} stores a new value in a place and returns the new
value (\dylanbook{chapter 7}). The place is either a variable name or a
call-like list.

When \meta{place} is a variable name, lexical or module, the variable is
given the value of \meta{new-value}. It is an error if there is no such
variable.

\begin{dylan}
(define counter 0)
(set! counter (+ counter 1))   ; => 1
counter                        ; => 1
(bind ((x 1))
  (set! x (* x 10))
  x)                           ; => 10
\end{dylan}

\noindent Some module variables cannot be assigned. A variable made by
\code{define-class} is read-only in both implementations. In the compiler
to C a variable holding a generic function cannot be assigned either, nor
can any variable of the \code{dylan} module; these are errors when
compiling. The Common Lisp implementation allows assigning a generic
function's variable, as the 1992 manual does.

\subsection{The extended form}\label{sec:control-extended-set}

When \meta{place} is a list \code{(\meta{name} \meta{arg} \ldots)}, it
stands for a location that the call \code{(\meta{name} \meta{arg} \ldots)}
would read, and \code{set!} stores into it by calling the setter function:

\begin{quote}
\code{(set! (\meta{name} \meta{arg} \ldots) \meta{new-value})}\\
means\\
\code{((setter \meta{name}) \meta{arg} \ldots{} \meta{new-value})}
\end{quote}

\noindent The new value comes \emph{last} among the setter's arguments.
Whatever the setter itself returns, \code{set!} returns the new value.

\begin{dylan}
(define v (vector 'a 'b 'c))
(set! (element v 1) 'sea)       ; => sea
v                               ; => #(a sea c)
((setter element) v 0 'x)       ; => x
v                               ; => #(x sea c)
\end{dylan}

\noindent The arguments and the new value are each evaluated once. The
compiler to C evaluates the arguments first, left to right, and the new
value last, the order of the call \code{set!} stands for; the Common Lisp
implementation evaluates the new value first. Slot accessors
(Chapter~\ref{ch:classes}) and \code{element} (Chapter~\ref{ch:collections}) are the
setters used most, and a program can define its own
(Section~\ref{sec:functions-setters}).

\section{Truth and if}\label{sec:control-if}

\code{\#f} is the only false value; every other object is true
(Section~\ref{sec:syntax-booleans}). A function described as ``returning true'' may
return any object other than \code{\#f}, not necessarily \code{\#t}.

\index{if@\code{if}}
\begin{quote}
\code{(if \meta{test} \meta{consequent} \meta{alternate})}
\end{quote}

\noindent \code{if} evaluates \meta{test}. If its value is true, \code{if}
evaluates \meta{consequent} and returns its values; otherwise it evaluates
\meta{alternate} and returns its values. Exactly three operands are
required: an \code{if} without an alternate, or with more than one form in
either branch, is an error in both implementations. Use \code{when} or
\code{unless} for a one-armed test, and \code{begin} to put several forms in
a branch.

\begin{dylan}
(if (< -11 0) 'negative 'non-negative)   ; => negative
(if 0 'yes 'no)                          ; => yes
(if '() 'yes 'no)                        ; => yes
(if #f 1 (values 4 5))                   ; => 4, 5
\end{dylan}

\index{not@\code{not}}
\code{(not \meta{thing})} is a function: it returns \code{\#t} if
\meta{thing} is \code{\#f}, and \code{\#f} otherwise.

\begin{dylan}
(not #f)                  ; => #t
(not 3)                   ; => #f
\end{dylan}

\section{Sequencing: begin}

\index{begin@\code{begin}}
\code{(begin \meta{form} \ldots)} evaluates the forms in order and returns
the values of the last one, or \code{\#f} if there are none. Many forms,
such as \code{bind}, \code{method} and \code{when}, already contain an
implicit \code{begin} around their bodies; an explicit one is needed where
a single form is expected, as in a branch of \code{if}.

\begin{dylan}
(if #t
    (begin (print "it was true")
           #t)
    #f)                  ; prints: "it was true"
                         ; => #t
(begin)                  ; => #f
\end{dylan}

\noindent A \code{begin} at top level makes its forms top-level forms too
(Section~\ref{sec:syntax-program}).

\section{Conditionals}\label{sec:control-conditionals}

\subsection{when and unless}

\index{when@\code{when}}\index{unless@\code{unless}}
\begin{quote}
\code{(when \meta{test} \meta{form} \ldots)}\\
\code{(unless \meta{test} \meta{form} \ldots)}
\end{quote}

\noindent \code{when} evaluates \meta{test}; if it is true, it evaluates
the forms and returns the values of the last. \code{unless} evaluates the
forms if \meta{test} is false. When the forms are not evaluated, or there
are none, the result is \code{\#f}.

\begin{dylan}
(when (> 3 2) 'big)      ; => big
(when (> 2 3) 'big)      ; => #f
(unless (> 2 3) 'a 'b)   ; => b
(when #t)                ; => #f
\end{dylan}

\subsection{and and or}

\index{and@\code{and}}\index{or@\code{or}}
\begin{quote}
\code{(and \meta{form} \ldots)}\\
\code{(or \meta{form} \ldots)}
\end{quote}

\noindent \code{and} evaluates the forms from left to right until one
returns \code{\#f} as its first value; then it returns \code{\#f} and
evaluates no more. If none does, it returns all the values of the last
form. \code{or} evaluates the forms from left to right until one returns a
true first value, and returns that value alone; if every form but the last
is false, it returns all the values of the last. With no forms, \code{and}
returns \code{\#t} and \code{or} returns \code{\#f}.

\begin{dylan}
(and 1 2 3)              ; => 3
(and 1 #f (error "never evaluated"))   ; => #f
(and 1 (values 2 3))     ; => 2, 3
(or #f 'x 'y)            ; => x
(or (values #f 3) 4)     ; => 4
(or #f (values 1 2))     ; => 1, 2
(and)                    ; => #t
(or)                     ; => #f
\end{dylan}

\subsection{cond}

\index{cond@\code{cond}}
\begin{quote}
\code{(cond (\meta{test} \meta{consequent} \ldots) \ldots)}
\end{quote}

\noindent \code{cond} evaluates the tests in order. At the first that is
true, it evaluates that clause's consequents and returns the values of the
last; later tests are not evaluated. A clause with no consequents returns
the value of its test. If no test is true, \code{cond} returns \code{\#f}.
By convention the last clause's test is \code{else:} or \code{\#t}; both are
simply true values.

\begin{dylan}
(define-method compare (new old)
  (cond ((< new old) "the new position is less")
        ((= new old) "the positions are equal")
        (else: "the new position is greater")))
(compare 1 2)            ; => "the new position is less"
(compare 3 2)            ; => "the new position is greater"
(cond (#f 'no) ((+ 1 2)))     ; => 3
(cond (#f 1))            ; => #f
\end{dylan}

\subsection{case}

\index{case@\code{case}}
\begin{quote}
\code{(case \meta{target} (\meta{match-list} \meta{consequent} \ldots) \ldots)}
\end{quote}

\noindent \code{case} evaluates \meta{target} and looks for its value in the
match lists, comparing with \code{id?}. The match lists are lists of
objects and are not evaluated. At the first clause whose list contains the
value, \code{case} evaluates that clause's consequents and returns the
values of the last, or \code{\#f} if it has none. A clause whose match list
is \code{else:} or \code{\#t} matches when no other does; it must be the last
clause. If nothing matches and there is no such clause, an error is
signaled.

\begin{dylan}
(define-method size-name (n)
  (case n
    ((0) 'none)
    ((1 2 3) 'few)
    (else: 'many)))
(size-name 2)            ; => few
(size-name 10)           ; => many
(case 'Red ((red green) 'warm-or-cool) (else: 'other))   ; => warm-or-cool
(case 9 ((3) 'three))    ; => error: no clause matches 9
\end{dylan}

\noindent Because the comparison is \code{id?}, matches are normally
numbers, characters, symbols, keywords, \code{\#t}, \code{\#f} or
\code{()}. Strings and lists are not suitable: two strings with the same
characters need not be the same object. The same match may not appear in
two clauses; the compiler to C reports that when compiling, and the Common
Lisp implementation uses the first.

\subsection{select}

\index{select@\code{select}}
\begin{quote}
\code{(select \meta{target} \meta{test} (\meta{match-list} \meta{consequent} \ldots) \ldots)}
\end{quote}

\noindent \code{select} is like \code{case}, but the matches are expressions,
evaluated, and the comparison is the function \meta{test}, called as
\code{(\meta{test} \meta{target-value} \meta{match})}. The clauses are tried
in order, and within a clause the matches from left to right; the matches
are evaluated only as they are reached, and once one matches, no more are
evaluated or tested (design note 34). \code{else:} and \code{\#t} work as in
\code{case}, and no match without them is an error.

\begin{dylan}
(define-method kind-of (x)
  (select x instance?
    ((<integer> <ratio>) "exact")
    ((<float>) "inexact")
    (else: "not a number")))
(kind-of 3)              ; => "exact"
(kind-of 2.5)            ; => "inexact"
(kind-of 'x)             ; => "not a number"
(select "b" = (("a") 1) (("b") 2))   ; => 2
(select 5 < ((3 4) 'small) ((6) 'medium) (else: 'large))   ; => medium
\end{dylan}

\noindent With a test of \code{instance?}, as in \code{kind-of},
\code{select} does what \code{typecase} does in Common Lisp.

\section{Iteration}\label{sec:control-loops}

The loops below are the language's iteration forms. Collections can also
be walked with functions such as \code{do} and \code{map}
(Chapters~\ref{ch:collections} and~\ref{ch:functional}), and any loop can be
written as a recursive method, since tail calls do not use stack
(Section~\ref{sec:functions-tail-calls}).

\subsection{for}

\index{for@\code{for}}
\begin{quote}
\code{(for ((\meta{var} \meta{init} \meta{step}) \ldots)\\
\phantom{(for }(\meta{end-test} \meta{result} \ldots)\\
\phantom{(for }\meta{body-form} \ldots)}
\end{quote}

\noindent \code{for} is the general loop. Each \code{(\meta{var}
\meta{init} \meta{step})} is a \term{variable clause}, and it is the only
kind of clause: each must have all three parts. The loop runs like this:

\begin{enumerate}
\item The inits are evaluated, in a scope that does not include any of the
  loop variables, and the variables are bound to their values.
\item The end test is evaluated. If it is true, the results are evaluated
  in order and \code{for} returns the values of the last, or \code{\#f} if
  there are none.
\item Otherwise the body forms are evaluated, for their effects.
\item The steps are evaluated, in a scope that includes the variables'
  current values, and the variables are bound afresh to the new values.
  The loop continues at step 2.
\end{enumerate}

\noindent The end-test clause may be the empty list \code{()}, meaning no
end test: the loop then runs until something exits from it
(Section~\ref{sec:control-exits}). Because all the steps are evaluated before any
variable is rebound, the variables are stepped in parallel.

\begin{dylan}
(define-method factorial ((n <integer>))
  (for ((i n (- i 1))
        (v 1 (* v i)))
       ((<= i 0) v)))
(factorial 5)            ; => 120
(for ((i 0 (+ i 1))
      (acc '() (cons i acc)))
     ((= i 4) acc))      ; => (3 2 1 0)
(for ((a 0 b) (b 1 (+ a b)) (i 0 (+ i 1)))
     ((= i 10) a))       ; => 55
\end{dylan}

\noindent Each pass binds the variables anew; it does not assign them. A
method made in the body therefore keeps the values of that pass:

\begin{dylan}
(define thunks
  (for ((i 0 (+ i 1))
        (acc '() (cons (method () i) acc)))
       ((= i 3) acc)))
(map (method (f) (f)) thunks)   ; => (2 1 0)
\end{dylan}

\subsection{for-each}

\index{for-each@\code{for-each}}
\begin{quote}
\code{(for-each ((\meta{var} \meta{collection}) \ldots)\\
\phantom{(for-each }(\meta{end-test} \meta{result} \ldots)\\
\phantom{(for-each }\meta{body-form} \ldots)}
\end{quote}

\noindent \code{for-each} iterates over the elements of one or more
collections at once. The collection expressions are evaluated first. Then,
on each pass, each \meta{var} is bound to the next element of its
collection; if any collection has no more elements, the loop ends and
returns \code{\#f}. Otherwise the end test is evaluated, with the variables
in scope; if it is true, the results are evaluated and the values of the
last returned (\code{\#f} if none). Otherwise the body forms are evaluated
and the next pass begins. As with \code{for}, the end-test clause may be
\code{()}, and each pass binds the variables anew.

\begin{dylan}
(define-method first-even ((s <sequence>))
  (for-each ((number s))
            ((even? number) number)))
(first-even '(1 3 5 4 7 9 10))   ; => 4
(first-even '(1 3))              ; => #f
(for-each ((c "abc")
           (i (range from: 0)))
          ()
  (print (list c i)))            ; prints: (#\a 0)
                                 ; prints: (#\b 1)
                                 ; prints: (#\c 2)
                                 ; => #f
\end{dylan}

\noindent The second example stops after three passes because the string
runs out, although the range \code{(range from: 0)} is unbounded. Any
collection may be used: lists, vectors, strings, ranges, tables
(whose elements are its values) and the program's own collection classes.
\code{for-each} uses the iteration protocol of Chapter~\ref{ch:collections}.

\subsection{dotimes}

\index{dotimes@\code{dotimes}}
\begin{quote}
\code{(dotimes (\meta{var} \meta{count} [\meta{result}]) \meta{form} \ldots)}
\end{quote}

\noindent \code{dotimes} evaluates \meta{count}, which should be an
integer, then evaluates the forms once for each integer from 0 up to but not
including the count, with \meta{var} bound to it. If the count is zero or
negative, the forms are not evaluated. Then \meta{result} is evaluated and
its values returned; without a result, \code{dotimes} returns \code{\#f}.
In both implementations the result form can see \meta{var}, bound to the
number of passes made (as in Common Lisp).

\begin{dylan}
(dotimes (i 3)
  (print i))             ; prints: 0
                         ; prints: 1
                         ; prints: 2
                         ; => #f
(dotimes (i 4 i))        ; => 4
(dotimes (i -2 'done))   ; => done
\end{dylan}

\subsection{while and until}

\index{while@\code{while}}\index{until@\code{until}}
\begin{quote}
\code{(while \meta{test} \meta{form} \ldots)}\\
\code{(until \meta{test} \meta{form} \ldots)}
\end{quote}

\noindent \code{while} evaluates \meta{test} before each pass and evaluates
the forms as long as it is true. \code{until} evaluates the forms as long as
\meta{test} is false. Both always return \code{\#f}.

\begin{dylan}
(define n 0)
(while (< n 3)
  (set! n (+ n 1)))      ; => #f
n                        ; => 3
(until (= n 0)
  (set! n (- n 1)))      ; => #f
n                        ; => 0
\end{dylan}

\section{Local methods: bind-methods}\label{sec:control-bind-methods}

\index{bind-methods@\code{bind-methods}}
\begin{quote}
\code{(bind-methods ((\meta{name} \meta{parameter-list} \meta{form} \ldots) \ldots)\\
\phantom{(bind-methods }\meta{body-form} \ldots)}
\end{quote}

\noindent \code{bind-methods} binds each \meta{name} to a method, made as
\code{method} makes one (Section~\ref{sec:functions-method}), and evaluates the
body forms, returning the values of the last. The scope of the names
includes the methods themselves, so the methods can call themselves and each
other. (\code{bind} cannot do that, since an init cannot see its own
variable.)

\begin{dylan}
(define-method root-mean-square ((s <sequence>))
  (bind-methods ((average (numbers)
                   (/ (reduce1 + numbers) (size numbers)))
                 (square (n) (* n n)))
    (sqrt (average (map square s)))))
(root-mean-square '(5 6 6 7 4))   ; => 5.692099788303083
\end{dylan}

\noindent The next example, Newton's method for square roots, has a method
that calls itself and two helpers:

\begin{dylan}
(define-method newtons-sqrt (x)
  (bind-methods ((sqrt1 (guess)
                   (if (close? guess)
                       guess
                       (sqrt1 (improve guess))))
                 (close? (guess)
                   (< (abs (- (* guess guess) x)) .0001))
                 (improve (guess)
                   (/ (+ guess (/ x guess)) 2)))
    (sqrt1 1)))
(newtons-sqrt 25.0)      ; => 5.000000000053722
\end{dylan}

\noindent The 1992 manual calls \code{(newtons-sqrt 25)}. With an integer
argument the arithmetic stays exact, and the guesses become ratios with
ever larger numerators and denominators: the Common Lisp implementation
returns an enormous ratio, and in the compiler to C, which has no bignums,
the integers overflow and an error is signaled. A float argument keeps the
computation in floating point.

Later versions of Dylan have a \code{local} form for this purpose; this
language has only \code{bind-methods}.

\section{Multiple values}\label{sec:control-values}

\index{values@\code{values}}\index{multiple values}
\begin{quote}
\code{(values \meta{object} \ldots)}
\end{quote}

\noindent An expression may return any number of values. The function
\code{values} returns its arguments as multiple values; \code{(values)}
returns none. \code{bind} (Section~\ref{sec:control-bind}) receives them.

\begin{dylan}
(define-method edges ((center <number>) (radius <number>))
  (values (- center radius) (+ center radius)))
(edges 100 2)            ; => 98, 102
(define-method opposite-edges ((center <number>) (radius <number>))
  (bind ((min max (edges center radius)))
    (values max min)))
(opposite-edges 100 2)   ; => 102, 98
(bind ((#rest nums (edges 100 2)))
  nums)                  ; => (98 102)
\end{dylan}

\noindent Multiple values pass through the forms that return the values
of a subform: \code{begin}, \code{bind}, \code{if}, \code{when},
\code{unless}, \code{cond}, \code{case}, \code{select}, the last form of
\code{and} and \code{or} as described above, the results of \code{for},
\code{for-each} and \code{dotimes}, \code{bind-exit}, \code{unwind-protect}'s
protected form, and calls of methods whose last form returns them. Anywhere
else, as an argument to a function, as the init of a one-variable binding,
as the value in a \code{set!} or a \code{define}, only the first value is
used, and an expression that returns no values supplies \code{\#f}.

\begin{dylan}
(begin 1 (values 2 3))   ; => 2, 3
(print (list (values 1 2) 3))   ; prints: (1 3)
(bind ((x (values))) x)  ; => #f
(values)                 ; => (no values)
\end{dylan}

\section{Non-local exits}\label{sec:control-exits}

\subsection{bind-exit}

\index{bind-exit@\code{bind-exit}}
\begin{quote}
\code{(bind-exit (\meta{exit-variable}) \meta{form} \ldots)}
\end{quote}

\noindent \code{bind-exit} binds \meta{exit-variable} to an \term{exit
procedure} and evaluates the forms. If the exit procedure is not called,
\code{bind-exit} returns the values of the last form, or \code{\#f} if
there are none. Calling the exit procedure, with any number of arguments,
leaves the \code{bind-exit} at once, and the arguments become its values.
This is a \term{non-local exit}.

\begin{dylan}
(bind-exit (return)
  1
  (return 2 3)
  4)                     ; => 2, 3
(bind-exit (return))     ; => #f
\end{dylan}

\noindent The exit procedure is an ordinary function: it can be passed to
other functions and stored. But it is valid only during the
\term{dynamic extent} of the \code{bind-exit} that made it, that is, until
that \code{bind-exit} returns or is itself exited. Calling it afterwards is
an error.

\begin{dylan}
(define-method first-even ((seq <sequence>))
  (bind-exit (exit)
    (do (method (item)
          (when (even? item)
            (exit item)))
        seq)
    #f))
(first-even '(1 3 5 4 7 9 10))   ; => 4
\end{dylan}

\noindent In the compiler to C a non-local exit is a jump and costs little;
\code{bind-exit} is the normal way to return early from a loop or a method.

\subsection{unwind-protect}

\index{unwind-protect@\code{unwind-protect}}
\begin{quote}
\code{(unwind-protect \meta{protected-form} \meta{cleanup-form} \ldots)}
\end{quote}

\noindent \code{unwind-protect} evaluates \meta{protected-form}, then the
cleanup forms. The cleanup forms are evaluated however the protected form
is left: normally, by a non-local exit through it, or by an error that a
handler outside it handles by exiting (Chapter~\ref{ch:conditions}). When the
protected form returns normally, \code{unwind-protect} returns its values;
the cleanup forms' values are discarded.

\begin{dylan}
(unwind-protect (values 1 2)
  (print 'cleanup))      ; prints: cleanup
                         ; => 1, 2
(bind-exit (k)
  (unwind-protect (k 5)
    (print 'cleanup)))   ; prints: cleanup
                         ; => 5
(handler-case
    (unwind-protect (error "boom")
      (print 'cleanup))
  ((<error>) 'caught))   ; prints: cleanup
                         ; => caught
\end{dylan}

\noindent A cleanup form may itself exit. If it does, the new exit replaces
the one in progress, and the remaining cleanup forms of that
\code{unwind-protect} are skipped:

\begin{dylan}
(bind-exit (k)
  (unwind-protect (k 1)
    (k 2)))              ; => 2
(bind-exit (one)
  (bind-exit (two)
    (unwind-protect (two 2)
      (one 1)))
  3)                     ; => 1
\end{dylan}

\noindent These are the rules design note 25 later gave for its
\code{block} form, which replaced \code{bind-exit} and
\code{unwind-protect} in later Dylan; this language keeps the two forms of
the 1992 manual.

\section{quote}

\code{(quote \meta{object})}, usually written \code{'\meta{object}}, returns
\meta{object} without evaluating it (Section~\ref{sec:syntax-quote}). It is a
special form of its own, listed here for completeness with the other
evaluation forms of \dylanbook{chapter 8}. \code{apply}, the last function
of that chapter, calls a function on a list of arguments and is described
with functions (Section~\ref{sec:functions-apply}).

\begin{dylan}
'+                       ; => +
(quote (+ 10 10))        ; => (+ 10 10)
\end{dylan}
